\setcounter{section}{19}






  \zwischenueberschrift{Übungsaufgaben}

\inputaufgabe
{}
{





Bestimme
\definitionsverweis {Idealerzeuger}{}{}
für die durch
\maabbeledisp {} { {\mathbb A}^{1}_{ K } } { { {\mathbb A}_{ K }^{  3   } }
} { t } { \left( t^2  , \,  t^3  , \,  t^4                 \right)
} {,}
gegebene
\definitionsverweis {monomiale Kurve}{}{.}





}
{} {}

\inputaufgabe
{}
{





Bestimme
\definitionsverweis {Idealerzeuger}{}{}
für die durch
\maabbeledisp {} { {\mathbb A}^{1}_{ K } } { { {\mathbb A}_{ K }^{  3   } }
} { t } { \left( t^4  , \,  t^5  , \,  t^6                 \right)
} {,}
gegebene
\definitionsverweis {monomiale Kurve}{}{.}





}
{} {}

\inputaufgabe
{}
{





Es sei

\mavergleichskette
{\vergleichskette
{R
}
{ \subseteq }{ S
}
{  }{
}
{    }{
}
{   }{
}
}
{}{}{}
eine
\definitionsverweis {ganze Erweiterung}{}{}
von
\definitionsverweis {Integritätsbereichen}{}{}
und sei

\mavergleichskette
{\vergleichskette
{F
}
{ \subseteq }{ R
}
{  }{
}
{    }{
}
{   }{
}
}
{}{}{}
ein
\definitionsverweis {multiplikatives System}{}{.}
Zeige, dass dann auch die zugehörige Erweiterung

\mavergleichskette
{\vergleichskette
{R_F
}
{ \subseteq }{ S_F
}
{  }{
}
{    }{
}
{   }{
}
}
{}{}{}
ganz ist.





}
{} {}

\inputaufgabegibtloesung
{}
{





Es seien $R$ und $S$
\definitionsverweis {Integritätsbereiche}{}{}
und sei

\mavergleichskette
{\vergleichskette
{R
}
{ \subseteq }{ S
}
{  }{
}
{    }{
}
{   }{
}
}
{}{}{}
eine
\definitionsverweis {ganze Ringerweiterung}{}{.}
Es sei

\mavergleichskette
{\vergleichskette
{f
}
{ \in }{ R
}
{  }{
}
{    }{
}
{   }{
}
}
{}{}{}
ein Element, das in $S$ eine Einheit ist. Zeige, dass $f$ dann schon in $R$ eine Einheit ist.





}
{} {}

\inputaufgabe
{}
{





Es sei

\mavergleichskette
{\vergleichskette
{ R
}
{ \subseteq }{ S
}
{  }{
}
{    }{
}
{   }{
}
}
{}{}{}
eine
\definitionsverweis {ganze Ringerweiterung}{}{}
und sei

\mavergleichskette
{\vergleichskette
{ f
}
{ \in }{ R
}
{  }{
}
{    }{
}
{   }{
}
}
{}{}{.}
Zeige: Wenn $f$, aufgefasst in $S$, eine
\definitionsverweis {Einheit}{}{}
ist, dann ist $f$ eine Einheit in $R$.





}
{} {}

\inputaufgabe
{}
{





Man gebe ein Beispiel einer
\definitionsverweis {ganzen Ringerweiterung}{}{}

\mavergleichskette
{\vergleichskette
{ R
}
{ \subseteq }{ S
}
{  }{
}
{    }{
}
{   }{
}
}
{}{}{,}
wo es einen
\definitionsverweis {Nichtnullteiler}{}{}

\mavergleichskette
{\vergleichskette
{ f
}
{ \in }{ R
}
{  }{
}
{    }{
}
{   }{
}
}
{}{}{}
gibt, der ein Nullteiler in $S$ wird.





}
{} {}

\inputaufgabe
{}
{





Es sei

\mavergleichskettedisp
{\vergleichskette
{P
}
{ =} { X^ 2-3X+7
}
{ } {
}
{ } {
}
{ } {
}
}
{}{}{}
und

\mavergleichskettedisp
{\vergleichskette
{Q
}
{ =} { Y^3-Y^2+4Y-5
}
{ } {
}
{ } {
}
{ } {
}
}
{}{}{.}
Begründe, dass die Ringerweiterung

\mavergleichskettedisp
{\vergleichskette
{\Z
}
{ \subseteq} { \Z[X,Y]/(P,Q)
}
{ } {
}
{ } {
}
{ } {
}
}
{}{}{}
\definitionsverweis {ganz}{}{}
ist und finde eine
\definitionsverweis {Ganzheitsgleichung}{}{}
für
\mathl{x+y}{} und für $xy$
\zusatzklammer {kleine Buchstaben bezeichnen die Restklassen der Variablen} {} {.}





}
{} {}

\inputaufgabe
{}
{





Es sei

\mavergleichskette
{\vergleichskette
{R
}
{ \subseteq }{ S
}
{  }{
}
{    }{
}
{   }{
}
}
{}{}{}
eine
\definitionsverweis {Ringerweiterung}{}{}
zwischen
\definitionsverweis {endlichen}{}{}
\definitionsverweis {kommutativen Ringen}{}{}
\mathkor {} {R} {und} {S} {.}
Zeige, dass eine
\definitionsverweis {ganze Ringerweiterung}{}{}
vorliegt.





}
{} {}

\inputaufgabe
{}
{





Es sei $R$ ein
\definitionsverweis {kommutativer Ring}{}{}
und

\mavergleichskettedisp
{\vergleichskette
{S
}
{ =} { R[X_1 , \ldots , X_n]/{\mathfrak a}
}
{ } {
}
{ } {
}
{ } {
}
}
{}{}{}
eine
\zusatzklammer {als Algebra} {} {}
\definitionsverweis {endlich erzeugte}{}{}
$R$-\definitionsverweis {Algebra}{}{,}
die
\definitionsverweis {ganz}{}{}
über $R$ sei. Zeige, dass $S$ ein
\definitionsverweis {endlich erzeugter}{}{}
$R$-\definitionsverweis {Modul}{}{}
ist.





}
{} {}

\inputaufgabe
{}
{





Es sei
\maabbdisp {\varphi} { R } { S
} {}
ein
\definitionsverweis {ganzer Ringhomomorphismus}{}{}
zwischen
\definitionsverweis {kommutativen Ringen}{}{}
und
\maabb {} { R } { R'
} {}
ein weiterer Ringhomomorphismus. Zeige, dass auch
\maabbeledisp {\varphi'} {R'} { R' \otimes_{ R } S
} { f } { f \otimes 1
} {,}
ganz ist.





}
{} {}

\inputaufgabe
{}
{





Zeige, dass der
\definitionsverweis {Ringhomomorphismus}{}{}
\maabbeledisp {} { K [X,Y]/(X^2+Y^2-1) } { K[U,V]/ \left(  U^2+V^2-1  \right)
} {(X,Y)} { \left(  U^2-V^2,2UV  \right)
} {,}
über jedem
\definitionsverweis {Körper}{}{}
der
\definitionsverweis {Charakteristik}{}{}
$\neq 2$
\definitionsverweis {ganz}{}{}
ist.





}
{} {}

\inputaufgabegibtloesung
{}
{





Es sei $K$ ein
\definitionsverweis {Körper}{}{}
und es sei eine
\definitionsverweis {endliche Ringerweiterung}{}{}
der Form

\mavergleichskettealigndrucklinks
{\vergleichskettealigndrucklinks
{ K[X]
}
{ \subseteq} { K[X] [Y]/ \left(  Y^n+ P_{n-1}(X)Y^{n-1} + \cdots + P_{2}(X)Y^{2} +P_1(X)Y+ P_0(X) \right)
}
{ =} { R
}
{ } {
}
{ } {
}
}
{}{}{}
mit

\mavergleichskette
{\vergleichskette
{P_i(X)
}
{ \in }{ K[X]
}
{  }{
}
{    }{
}
{   }{
}
}
{}{}{}
gegeben, wobei der Erweiterungsring
\definitionsverweis {integer}{}{}
sei. Es sei $k$ derart, dass der Grad von $P_i$ höchstens
\mathl{k(n-i)}{} ist für alle

\mavergleichskettedisp
{\vergleichskette
{ i
}
{ =} { 0,1 , \ldots , n-1
}
{ } {
}
{ } {
}
{ } {
}
}
{}{}{.}
Zeige, dass dann $YX^{-k}$ eine Ganzheitsgleichung vom Grad $n$ über $K[ X^{-1} ]$ erfüllt, und dass die zugehörige Erweiterungsalgebra
\mathl{K[ X^{-1},  YX^{-k} ]}{}  den gleichen Quotientenkörper wie $R$  besitzt.





}
{} {}

\inputaufgabe
{}
{





\aufzaehlungzwei {Es sei $R$ ein
\definitionsverweis {Integritätsbereich}{}{.}
Zeige, dass $R$
\definitionsverweis {ganz-abgeschlossen}{}{}
im
\definitionsverweis {Polynomring}{}{}

\mathl{R[X]}{} ist.
} { Man gebe ein Beispiel für einen kommutativen Ring $R$, der im Polynomring nicht ganz-abgeschlossen ist.
}





}
{} {}





  \zwischenueberschrift{Aufgaben zum Abgeben}

\inputaufgabe
{4}
{





Es sei

\mavergleichskette
{\vergleichskette
{ M
}
{ \subseteq }{ \N
}
{  }{
}
{    }{
}
{   }{
}
}
{}{}{}
das durch
\mathl{3,5,7}{} erzeugte numerische Untermonoid. Bestimme eine Restklassendarstellung des zugehörigen Monoidringes.





}
{} {}

\inputaufgabe
{3}
{





Es seien
\mathl{R,S,T}{}
\definitionsverweis {kommutative Ringe}{}{}
und seien
\mathl{\varphi:R \rightarrow S}{} und
\mathl{\psi:S \rightarrow T}{} Ringhomomorphismen derart, dass $S$
\definitionsverweis {ganz}{}{} über $R$ und $T$ ganz über $S$ ist. Zeige, dass dann auch $T$ ganz über $R$ ist.





}
{} {\zusatzklammer {Vergleiche

Aufgabe 10.26
} {} {.}}
